\section{速度与任务适配性}

\subsection{速度不是决定性因素}

Claude Code 公认速度更快，但作者指出：如果一个代理快速完成任务却留下需要 10 分钟调试的 bug，另一个虽慢却直接交付可用结果，后者的额外时间完全值得。

\begin{warningbox}{速度陷阱}
不要被"我的代理更快"这类说法误导。编程代理是长期使用的工具，\textbf{端到端的完成质量}比单次任务的执行速度更重要。
\end{warningbox}

\subsection{任务类型决定胜负}

两款代理在不同编程领域的表现差异很大：

\begin{itemize}
    \item 在 AI Engineering 任务中，一方可能占优
    \item 在 Web 开发任务中，另一方可能反超
    \item 低级编程（low-level programming）领域尚无定论
\end{itemize}

目前缺乏系统性的跨任务对比研究，且由于模型和代理每隔几个月就会大幅更新，这类研究也很难长期有效。

\subsection{本章小结}

速度优势不等于体验优势，任务类型对两款代理的表现影响极大。理想做法是在自己的实际编程领域中小规模测试两者。

